\scanpage{14}

17) Soit $C\in G'(\mathbb R)$ \deleted{\unsure{$\varepsilon$}}, avec $C^2=\unsure{y|\gamma|\eta}$ $(=i(-1))$, définissant une polarisation de $\mathcal M=\operatorname{Rep}(G)$, -- d'un $G$ réductif !

Alors $\overbrace{\operatorname{Centr}_{G'}(C)}^{K}(\mathbb R)$ est un \uline{compact} maximal de $G'(\mathbb R)$. En effet, on \deleted{\unsure{voit tout de suite que}} \deleted{\unsure{$\varphi$}} aura \deleted{$G'\subset\operatorname{Aut}_{\mathbb R}(M^0)\times\operatorname{Aut}_{\mathbb R}(M^1)$, et} $G'\hookrightarrow\prod_\alpha\operatorname{Aut}_{\mathbb R}(M_\alpha)$, par des repr. convenables (au nb de 2 si on veut) de $G'$, \deleted{de} avec $y$ y opérant par $\pm1$. Si $\varphi_\alpha$ est une forme de $C$-polarisation sur $M_\alpha$, alors les $g\in\operatorname{Centr}_{G'}(C)$ laissent invariant $\varphi_\alpha^C$, donc

\[
\deleted{\unsure{\operatorname{Centr}|\operatorname{Cent}}}\quad K\subset\prod_\alpha\underbrace{\operatorname{Aut}_{\mathbb R}(M_\alpha,\varphi_\alpha^C)}_{\substack{\text{groupes réels}\\\text{compacts}}}\ ,
\]

donc $K$ est \inserted{en pts réels} compact. Prouvons que \deleted{il est} $K(\mathbb R)$ est compact \uline{maximal}. Soit \deleted{donc} $\underline{K'}\subset G'(\mathbb R)$, $K'\supset K(\mathbb R)$. On sait que $\underline{K'}$ est algébrique, $\underline{K'}=K'(\mathbb R)$, $K'$ ss-groupe algébrique de $G'$ contenant $C$. Quitte à remplacer $G'$ par $K'$, on peut supposer $G'$ compact, et il faut prouver qu'alors $C$ est central.

Passons le d'abord dans $G'/\{1,y\}$, i.e. supposons d'abord $y=1$. Soit une repr. fidèle de $G'$ dans un $V$, \deleted{\unsure{avec forme de pol. sur} $V$} \deleted{Alors $C$ \unsure{définit}} \inserted{avec forme de $C$-pol. $\varphi$ sur $V$.} $V=V^++V^-$

\deleted{($C\mid V^+=\mathrm{id}$, $C\mid V^-=-\mathrm{id}$) et} avec $(\varphi\mid V^+)$ \deleted{est} $>0$, $\underset{\substack{\Vert\\-\varphi^-}}{(\varphi\mid V^-)}<0$.

\deleted{\unsure{$\varepsilon$}} $G'$ étant un sous-groupe \unsure{Cpt|cpt} de $O(\varphi)$, on peut par conjugaison \inserted{par $u\in O(\varphi)(\mathbb R)$} l'envoyer dans le compact maximal $O(V^+)\times O(V^-)$ \inserted{de $O(\varphi)$}, donc on peut supposer que $G\subset O(V^+)\times O(V^-)$. Donc $C\in G(\mathbb R)$ est de la forme $C^+\oplus C^-$, \deleted{\unsure{et alors}} Comme $\varphi^+(x,C^+y)$ et

% archive sheet number: 13
